% Standalone section for incorporation into the main paper.
% All arguments below are independent of Dirichlet GRH.
\section{The whole corank pattern of a random alternating pencil}

Write $\operatorname{Alt}_n(q)$ for the vector space of alternating bilinear
forms on $\mathbb F_q^n$, including zero diagonal in characteristic two. Put
\[
 g_j=\prod_{i=1}^j(1-q^{-i}),\qquad
 g_\infty=\prod_{i\ge1}(1-q^{-i}),\qquad
 \beta_q=\prod_{i\ge1}(1-q^{-(2i-1)}).
\]
All limits in this section keep $q$ and all coefficient directions fixed,
and send $n$ to infinity through even integers.

\begin{theorem}[Joint coranks at every fixed projective direction]
\label{thm:pencil-coranks}
Let $A_1,\ldots,A_k$ be independent uniform elements of
$\operatorname{Alt}_n(q)$. Choose nonzero vectors
$a_1,\ldots,a_s\in\mathbb F_q^k$, no two proportional, and put
$B_i=\sum_j a_{ij}A_j$. Then
\[
 \bigl(\dim\ker B_1,\ldots,\dim\ker B_s\bigr)
 \ \longrightarrow\ (R_1,\ldots,R_s)
\]
in total variation, where the $R_i$ are independent and have distribution
\[
 \Pr(R_i=2h)=\frac{\beta_q q^{-h(2h-1)}}{g_{2h}},
 \qquad h=0,1,2,\ldots.
\]
In particular the nonsingularity indicators are asymptotically independent
Bernoulli variables of parameter $\beta_q$.
\end{theorem}

The independence assertion allows arbitrary linear dependencies among the
coefficient vectors; only distinct projective directions are required.

\begin{lemma}[Strict tensor dimension surplus]
\label{lem:tensor-surplus}
Let $a_i$ be pairwise nonproportional nonzero vectors in a vector space $E$
over any field, and let $W_i$ be finite-dimensional subspaces of another
vector space. Set
\[
 U=\sum_iW_i,\qquad T=\sum_i a_i\otimes W_i\subseteq E\otimes U.
\]
Then $\dim T\ge\dim U$, with equality if and only if the sum of the $W_i$
is direct.
\end{lemma}

\begin{proof}
Pass to an infinite extension field; this preserves dimensions and directness.
Choose $f\in E^*$ nonzero on every $a_i$. Contraction by $f$ maps $T$ onto
$U$, proving the inequality. If the dimensions are equal this contraction
is an isomorphism. Choose $g\in E^*$ so that all ratios
$g(a_i)/f(a_i)$ are distinct; finitely many proper hyperplanes are excluded.
The operator
\[
 (g\otimes1)\circ(f\otimes1)^{-1}:U\longrightarrow U
\]
acts on $W_i$ by the scalar $g(a_i)/f(a_i)$. Eigenspaces with distinct
eigenvalues have direct sum. Conversely a direct sum of the $W_i$ makes the
displayed tensor sum direct, of dimension $\sum_i\dim W_i=\dim U$.
\end{proof}

For $r,m\ge0$ define the number of injective maps
\[
 I_m(r)=\prod_{j=0}^{m-1}(q^r-q^j).
\]
The empty product is one, and $I_m(r)=0$ when $m>r$.

\begin{lemma}[All mixed injection moments]
\label{lem:mixed-injection-moments}
For fixed nonnegative integers $m_1,\ldots,m_s$ and $r_i=\dim\ker B_i$,
\[
 \lim_n\mathbb E\prod_i I_{m_i}(r_i)
   =\prod_i q^{m_i(m_i+1)/2}.
\]
Consequently, for fixed nonnegative integers $\ell_i$,
\[
 \lim_n\mathbb E\prod_i q^{\ell_i r_i}
 =\prod_i M_{\ell_i},\qquad
 M_\ell=\prod_{j=1}^{\ell}(1+q^j).
\]
\end{lemma}

\begin{proof}
Expand the injection counts over tuples of injective maps
$f_i:\mathbb F_q^{m_i}\to\mathbb F_q^n$, and let $W_i$ be their images.
Write $M=\sum_i m_i$.

First suppose the $W_i$ have direct sum. The number of such ordered tuples
of injections is $I_M(n)$. In a basis adapted to their sum, the constraints
$W_i\subseteq\ker B_i$ have total codimension
\[
 M(n-M)+\sum_i\binom{m_i}{2}+2\sum_{i<j}m_im_j
 =nM-\sum_i\frac{m_i(m_i+1)}2.
\]
Indeed, each block joining $W_i$ to the complementary space imposes one
linear condition per entry; each alternating block within $W_i$ imposes
one; and a block joining $W_i$ to $W_j$ imposes two independent conditions,
because $a_i$ and $a_j$ are nonproportional. All underlying upper-triangular
entries of the independent $A_j$ are uniform. The contribution of direct
tuples is therefore
\[
 I_M(n)q^{-nM+\sum_i m_i(m_i+1)/2}
 \longrightarrow q^{\sum_i m_i(m_i+1)/2}.
\]

For a nondirect tuple put $u=\dim\sum_iW_i<M$ and
$v=\dim\sum_i a_i\otimes W_i$. Lemma~\ref{lem:tensor-surplus} gives
$v\ge u+1$. For each of the $n-u$ complementary basis vectors, the
kernel conditions impose $v$ independent linear equations on the
$k\times u$ cross-block entries. These entries for different complementary
vectors are independent. Thus the probability of all kernel conditions
is at most $q^{-v(n-u)}$. The number of tuples spanning a space of dimension
$u$ is at most $C_{q,M}q^{nu}$: choose the $u$-space, then the bounded-size
coordinate matrices of the injections within it. Their contribution is
at most $C_{q,M}q^{-n+v u}$, which tends to zero. There are only finitely many
possible $u,v$ for the fixed moment tuple.

Finally, counting maps by their kernels gives
\[
 q^{\ell r}=\sum_{j=0}^{\ell}\qbinom{\ell}{j} I_j(r).
\]
Apply the first part and the finite $q$-binomial identity
\[
 \sum_{j=0}^{\ell}\qbinom{\ell}{j} q^{j(j+1)/2}
 =\prod_{j=1}^{\ell}(1+q^j).
\]
\end{proof}

\begin{proof}[Proof of Theorem~\ref{thm:pencil-coranks}]
We give the moment-determination step explicitly, since ordinary moment
convergence by itself would not justify this conclusion. Put
\[
 B_q=\prod_{j\ge1}(1-q^{-2j}),\qquad
 C_q^+=\prod_{j\ge1}(1+q^{-j}).
\]
Both are finite positive constants. For fixed $h\ge0$ and $m\ge h$, define
the polynomial in $q^r$
\[
 L_{m,h}(r)=\prod_{\substack{0\le j\le m\\j\ne h}}
            \frac{q^r-q^{2j}}{q^{2h}-q^{2j}}.
\]
Every alternating form in even dimension has even nullity, over every
finite field. On the even values $0,2,\ldots,2m$, this polynomial equals
$\mathbf1_{\{r=2h\}}$. If $r\ge2m+2$ is even, bounding the numerator
factors by $q^r$ and the denominator factors separately above and below
$h$ gives
\[
 |L_{m,h}(r)|
 \le B_q^{-2}q^{mr-m(m+1)-h(h-1)}.
\]
Consequently, at every even $r\ge0$,
\[
 \left|L_{m,h}(r)-\mathbf1_{\{r=2h\}}\right|
 \le E_{m,h}(r):=
 B_q^{-2}q^{-h(h-1)}q^{(m+1)r-(m+1)(m+2)}.
\]
By Lemma~\ref{lem:mixed-injection-moments}, all fixed mixed moments of
these error majorants factor in the limit. Moreover,
\[
 \lim_n\mathbb E E_{m,h}(r_i)
 \le B_q^{-2}C_q^+q^{-h(h-1)}q^{-(m+1)(m+2)/2}
 =:\varepsilon_{m,h}\longrightarrow0,
\]
since $M_{m+1}\le C_q^+q^{(m+1)(m+2)/2}$.
For fixed $h_1,\ldots,h_s$, the difference between the product of the
indicators and the product of the polynomials is bounded by
$\prod_i(1+E_{m,h_i}(r_i))-1$. Its limiting expected value is at most
$\prod_i(1+\varepsilon_{m,h_i})-1$, tending to zero with $m$.
The mixed polynomial expectations factor in the limit, again by the
mixed moment lemma. Thus each joint point probability converges to the
product of its marginal limits.

For completeness, choosing the radical and then a nondegenerate form on
the quotient gives, for even $r\le n$,
\[
 \Pr(\dim\ker A_1=r)
 =\qbinom{n}{r}
 q^{\binom{n-r}{2}-\binom n2}
 \prod_{j=1}^{(n-r)/2}(1-q^{-(2j-1)}).
\]
Its limit is $\beta_q q^{-r(r-1)/2}/g_r$. The first moment bound
$\mathbb E q^{r_i}\to q+1$ makes the marginals tight, so these probabilities
sum to one. Joint tightness follows by a union bound. Pointwise
convergence on the countable state space, together with tightness, gives
total-variation convergence.
\end{proof}

\begin{corollary}[A random finite-dimensional space of alternating forms]
\label{cor:binomial-pencil}
Let $U_n$ be a uniformly chosen $k$-dimensional subspace of
$\operatorname{Alt}_n(q)$, where $k$ is fixed. Put
$m_k=(q^k-1)/(q-1)$, and let $Z(U_n)$ be the number of its projective points
represented by nonsingular forms. Then
\[
 Z(U_n)\longrightarrow\operatorname{Bin}(m_k,\beta_q)
\]
in total variation. More strongly, with respect to a uniformly chosen
ordered basis of $U_n$, the entire projective corank pattern converges to
independent copies of the distribution in Theorem~\ref{thm:pencil-coranks}.
In particular,
\[
 \Pr(\text{every nonzero form of }U_n\text{ is nonsingular})
 \longrightarrow\beta_q^{m_k}.
\]
\end{corollary}

\begin{proof}
Independent uniform $A_1,\ldots,A_k$ are linearly independent with
probability $\prod_{j=0}^{k-1}(1-q^{j-\binom n2})\to1$. Conditioning on
that event gives a uniform ordered basis of a uniform $k$-space. Apply
the theorem to one representative of every point of
$\mathbb P^{k-1}(\mathbb F_q)$. The case $k=0$ is interpreted with zero
projective points.
\end{proof}

\paragraph{Relation to prior work.}
The tensor-dimension mechanism has a closely related unrestricted-matrix
antecedent in the November 2026 Omega Institute proof of the fixed-field
column limit for MRD-code densities \cite{OmegaMRD2026}. That work uses mixed kernel moments
and finite Mobius inversion to establish joint zero-nullity probabilities.
Here the matrices are alternating, every nullity has fixed parity, and
the conclusion determines the entire projective corank pattern. The
explicit interpolation argument above is included to establish this
stronger distributional statement without importing a moment-uniqueness
theorem. This comparison is a targeted prior-art check, not a claim of
exhaustive priority verification.
